
%修改 proof 证明环境
\makeatletter
\renewcommand\proofname{证明}
\renewenvironment{proof}[1][\proofname]{\par
\pushQED{\qed}%
\normalfont \topsep6\p@\@plus6\p@ \labelsep1em\relax
\trivlist
\item[\hskip\labelsep\indent
\bfseries #1]\ignorespaces
}{%
\popQED\endtrivlist\@endpefalse
}
\makeatother

% 例题环境
\newcounter{example}[section]
\renewcommand{\theexample}{\thesection.\arabic{example}}

\newenvironment{example}[1][]{\refstepcounter{example} \textbf{例 \theexample  \ #1} \hspace{0.5em} \itshape}{\hspace{\stretch{1}} \rule{1ex}{1ex}}

% 习题环境
\newcounter{exercise}[section]
\renewcommand{\theexercise}{\thesection.\arabic{exercise}}

\newenvironment{exercise}[1][]{\refstepcounter{exercise} \textbf{题 \theexercise  \ #1} \hspace{0.5em}}{\hspace{\stretch{1}} \rule{1ex}{1ex}}

\newcommand{\exerciseFrom}[1][]{\textbf{题目出处}\hspace{1em}#1}
\newcommand{\exerciseSolvedDate}[1][]{\textbf{解答日期}\hspace{1em}#1}
\newenvironment{exerciseAdditional}{\textbf{补充说明}\hspace{1em}}{}

\newtheorem{definition}{定义}[section]
\newtheorem{property}{性质}[section]
\newtheorem{theorem}{定理}[section]
\newtheorem{inference}{推论}[section]
\newtheorem{axiom}{公理}[section]
\newtheorem{lemma}{引理}[section]

\usepackage[most]{tcolorbox}
\usepackage{xcolor} % 确保颜色可用
\tcbuselibrary{theorems, skins} 
% 引入fontawesome5以使用图标（可选，但推荐增加视觉效果）
\usepackage{fontawesome5}
\usepackage[framemethod=TikZ]{mdframed} % 用于创建带边框的环境

% 使用 newmdtheoremenv 命令来定义一个新的mdframed定理环境
\newmdtheoremenv[
    linecolor=blue!75!black,  % 边框颜色为蓝色
    linewidth=1pt,            % 边框粗细为 1pt
    frametitlefont=\bfseries\large, % 标题字体加粗
    backgroundcolor=blue!5!white, % 背景色为浅蓝色
    roundcorner=5pt,       % 圆角设计
    skipabove=1em,         % 环境上方留白
    skipbelow=1em,         % 环境下方留白
    innertopmargin=1.5em,
    innerbottommargin=1em,
    innerleftmargin=1em,
    innerrightmargin=1em,
]{principle}{原理}[section]

%\newtheorem{principle}{原理}[section]
%\newtheorem{exercise}{题目}[section]
\newtheorem{topic}{问题}[section]
\newtheorem{statement}{命题}[section]
% \newtheorem{example}{例}[section]

% 使公式编号与章节关联，命令由 amsmath 宏包提供
\numberwithin{equation}{section}

% 配合 \autoref 命令使引用不只引用编号，也能引用环境命名，如: 定理 3.2.5，来自 hyperref 宏包。
%\newcommand\equationautorefname{式}
%\newcommand\footnoteautorefname{脚注}%
%\newcommand\itemautorefname{项}
\def\figureautorefname{图}
\def\tableautorefname{表}
\def\chapterautorefname{章}
\def\sectionautorefname{节}
\def\subsectionautorefname{小节}
\def\appendixautorefname{附录}
\def\propertyautorefname{性质}
\def\theoremautorefname{定理}
\def\definitionautorefname{定义}
\def\inferenceautorefname{推论}
\def\axiomautorefname{公理}
\def\lemmaautorefname{引理}
\def\principleautorefname{原理}
\def\exerciseautorefname{题目}
\def\topicautorefname{问题}
\def\statementautorefname{命题}
\def\exampleautorefname{例}
\def\equationautorefname{式}

\newcommand{\pair}[2]{\langle #1, #2 \rangle}



% --- 定义核心思想 tcolorbox 环境 ---
\newtcolorbox{coreidea}[1]{
    %breakable,          % 允许跨页
    enhanced,           % 启用增强效果
    colback=white,      % 设置主体文字颜色
    colframe=gray,      % 设置主体框体边框的颜色
    colbacktitle=white, % 设置标题区域的背景颜色
    coltitle=gray,     % 设置标题文本的颜色
    attach boxed title to top left = {yshift = -4mm, xshift = 8mm}, % 将标题定位到框的上边缘
    boxsep=4mm, % 个方向空隙
    bottom=-1mm, % 底部抵消一部分
    fonttitle=\bfseries\large, % 标题字体样式
    fontupper=\sffamily,
    title={\faIcon{lightbulb}\,\,#1},
    boxrule=1pt,        % 边框粗细
    arc=3pt,            % 圆角
    drop fuzzy shadow,  % 添加阴影
    boxed title style={
      frame code={
        \draw[tcb fill frame,draw=gray](
        [xshift=-4mm]frame.west)--(frame.north west)
        --(frame.north east) --([xshift=4mm]frame.east)
        --(frame.south east)--(frame.south west) --([xshift=-4mm]frame.west);
      },
      interior code={
        \draw[tcb fill interior,draw=gray](
        [xshift=-2mm]interior.west)--(interior.north west)
        -- (interior.north east) --([xshift=2mm]interior.east)
        --(interior.south east)--(interior.south west) --([xshift=-2mm]interior.west);
      }
    }
}

% --- 使用 tcolorbox 定义 principle 环境 ---
\tcbset{
    principlebox/.style={
        enhanced,
        colback=white,% 背景色
        colframe=blue!75!black,% 边框颜色
        coltitle=black,% 标题颜色
        boxrule=1pt,           % 边框粗细
        arc=5pt,              % 设置圆角半径
        % *** 关键设置 1: 将标题附加到框的外部左上角 ***
        attach boxed title to top left={
          xshift=10pt,
          yshift=-7pt,
        }, 
        % 关键设置 2: 标题区域的样式
        boxed title style={
            colback=white,
            arc=3pt,
            boxsep=3pt, % 标题框内部的留白
        },
        % *** 关键设置 3: 确保正文内容与标题区域有足够垂直距离，强制另起一行 ***
        top=1em, 
        fonttitle=\bfseries, % 标题字体
        fontupper=\itshape,
    }
}

% 使用 newtcbtheorem 定义带有编号的定理环境
% 参数顺序：{计数器名称}{显示的标题文字}{tcolorbox样式名称}{标签前缀}
\newtcbtheorem[number within=section]{myprinciple}{原理}{principlebox}{myprinciple}
% --------------------------------------------------

%%% Local Variables:
%%% mode: latex
%%% TeX-master: "elementary-math-note"
%%% End:
